% GENERATED FILE. Edit canonical lessons, then run npm run generate:books.
% canonical-source-hash: fnv1a64:521c2e984256e7f8
% canonical-lessons: ML-C222-bhumisastram, ML-C222-sahityam, ML-C222-vyakaranam, ML-C222-barbar, ML-C222-tayyalkkaran, ML-R222-first-pass-words-from-chapters-217-219, ML-R222-second-pass-words-from-chapters-220-222

\chapter{Geography, Literature, Grammar, Barber, Tailor}
\label{ch:ml-geography-literature-grammar-barber-tailor}
\hlchaptermodality{222}

\begin{chapteropening}
\textbf{By the end of this chapter:} \emph{I can say geography, literature, grammar, a barber and a tailor.}

Complete the last lesson of chapter 222: I can say geography, literature, grammar, a barber and a tailor.
\end{chapteropening}

\section[\texorpdfstring{\ml{ഭൂമിശാസ്ത്രം} (bhūmiśāstraṁ)}{ഭൂമിശാസ്ത്രം (bhūmiśāstraṁ)}]{\ml{ഭൂമിശാസ്ത്രം} (bhūmiśāstraṁ) — geography}

\label{lesson:ML-C222-bhumisastram}



\begin{quote}
Before the new one: say the Malayalam for a classmate, then the Malayalam for education.
\end{quote}

\subsection*{You'll want to know: \ml{ഭൂമിശാസ്ത്രം}}
\textbf{\ml{ഭൂമിശാസ്ത്രം}} — \emph{bhūmiśāstraṁ} — \textquotedblleft{}geography\textquotedblright{}.

\begin{grammarlens}[title={What you've built}]
One more word for this chapter.
\end{grammarlens}

\subsection*{Your turn}
\begin{itemize}
\raggedright
  \item \emph{Say it:} \emph{bhūmiśāstraṁ}
  \item \emph{Say it:} \emph{bhūmiśāstraṁ}, once more
  \item \emph{Say it:} say \emph{bhūmiśāstraṁ}
  \item \emph{Recall:} say the Malayalam for a classmate, then the Malayalam for education, then say \emph{bhūmiśāstraṁ} again
\end{itemize}

\begin{tcolorbox}[breakable,colback=teal!4,colframe=teal!35!black,title={Before you move on}]
What does \ml{ഭൂമിശാസ്ത്രം} mean? (\textquotedblleft{}Geography\textquotedblright{}.) Say it once more.
\end{tcolorbox}
\section[\texorpdfstring{\ml{സാഹിത്യം} (sāhityaṁ)}{സാഹിത്യം (sāhityaṁ)}]{\ml{സാഹിത്യം} (sāhityaṁ) — literature}

\label{lesson:ML-C222-sahityam}



\begin{quote}
Before the new one: say the Malayalam for education, then the Malayalam for geography.
\end{quote}

\subsection*{You'll want to know: \ml{സാഹിത്യം}}
\textbf{\ml{സാഹിത്യം}} — \emph{sāhityaṁ} — \textquotedblleft{}literature\textquotedblright{}.

\begin{grammarlens}[title={What you've built}]
One more word for this chapter.
\end{grammarlens}

\subsection*{Your turn}
\begin{itemize}
\raggedright
  \item \emph{Say it:} \emph{sāhityaṁ}
  \item \emph{Say it:} \emph{sāhityaṁ}, once more
  \item \emph{Say it:} say \emph{sāhityaṁ}
  \item \emph{Recall:} say the Malayalam for education, then the Malayalam for geography, then say \emph{sāhityaṁ} again
\end{itemize}

\begin{tcolorbox}[breakable,colback=teal!4,colframe=teal!35!black,title={Before you move on}]
What does \ml{സാഹിത്യം} mean? (\textquotedblleft{}Literature\textquotedblright{}.) Say it once more.
\end{tcolorbox}
\section[\texorpdfstring{\ml{വ്യാകരണം} (vyākaraṇaṁ)}{വ്യാകരണം (vyākaraṇaṁ)}]{\ml{വ്യാകരണം} (vyākaraṇaṁ) — grammar}

\label{lesson:ML-C222-vyakaranam}



\begin{quote}
Before the new one: say the Malayalam for geography, then the Malayalam for literature.
\end{quote}

\subsection*{You'll want to know: \ml{വ്യാകരണം}}
\textbf{\ml{വ്യാകരണം}} — \emph{vyākaraṇaṁ} — \textquotedblleft{}grammar\textquotedblright{}.

\begin{grammarlens}[title={What you've built}]
One more word for this chapter.
\end{grammarlens}

\subsection*{Your turn}
\begin{itemize}
\raggedright
  \item \emph{Say it:} \emph{vyākaraṇaṁ}
  \item \emph{Say it:} \emph{vyākaraṇaṁ}, once more
  \item \emph{Say it:} say \emph{vyākaraṇaṁ}
  \item \emph{Recall:} say the Malayalam for geography, then the Malayalam for literature, then say \emph{vyākaraṇaṁ} again
\end{itemize}

\begin{tcolorbox}[breakable,colback=teal!4,colframe=teal!35!black,title={Before you move on}]
What does \ml{വ്യാകരണം} mean? (\textquotedblleft{}Grammar\textquotedblright{}.) Say it once more.
\end{tcolorbox}
\section[\texorpdfstring{\ml{ബാർബർ} (bārbar)}{ബാർബർ (bārbar)}]{\ml{ബാർബർ} (bārbar) — a barber}

\label{lesson:ML-C222-barbar}



\begin{quote}
Before the new one: say the Malayalam for literature, then the Malayalam for grammar.
\end{quote}

\subsection*{You'll want to know: \ml{ബാർബർ}}
\textbf{\ml{ബാർബർ}} — \emph{bārbar} — \textquotedblleft{}a barber\textquotedblright{}.

\begin{grammarlens}[title={What you've built}]
One more word for this chapter.
\end{grammarlens}

\subsection*{Your turn}
\begin{itemize}
\raggedright
  \item \emph{Say it:} \emph{bārbar}
  \item \emph{Say it:} \emph{bārbar}, once more
  \item \emph{Say it:} say \emph{bārbar}
  \item \emph{Recall:} say the Malayalam for literature, then the Malayalam for grammar, then say \emph{bārbar} again
\end{itemize}

\begin{tcolorbox}[breakable,colback=teal!4,colframe=teal!35!black,title={Before you move on}]
What does \ml{ബാർബർ} mean? (\textquotedblleft{}A barber\textquotedblright{}.) Say it once more.
\end{tcolorbox}
\section[\texorpdfstring{\ml{തയ്യൽക്കാരൻ} (tayyalkkāran)}{തയ്യൽക്കാരൻ (tayyalkkāran)}]{\ml{തയ്യൽക്കാരൻ} (tayyalkkāran) — a tailor}

\label{lesson:ML-C222-tayyalkkaran}



\begin{quote}
Before the new one: say the Malayalam for grammar, then the Malayalam for a barber.
\end{quote}

\subsection*{You'll want to know: \ml{തയ്യൽക്കാരൻ}}
\textbf{\ml{തയ്യൽക്കാരൻ}} — \emph{tayyalkkāran} — \textquotedblleft{}a tailor\textquotedblright{}.

\begin{grammarlens}[title={What you've built}]
One more word for this chapter.
\end{grammarlens}

\subsection*{Your turn}
\begin{itemize}
\raggedright
  \item \emph{Say it:} \emph{tayyalkkāran}
  \item \emph{Say it:} \emph{tayyalkkāran}, once more
  \item \emph{Say it:} say \emph{tayyalkkāran}
  \item \emph{Recall:} say the Malayalam for grammar, then the Malayalam for a barber, then say \emph{tayyalkkāran} again
\end{itemize}

\begin{tcolorbox}[breakable,colback=teal!4,colframe=teal!35!black,title={Before you move on}]
What does \ml{തയ്യൽക്കാരൻ} mean? (\textquotedblleft{}A tailor\textquotedblright{}.) Say it once more.
\end{tcolorbox}
\section[(dialogue)]{Words from chapters 217-219 — a first pass}

\label{lesson:ML-R222-first-pass-words-from-chapters-217-219}



\begin{quote}
Say the words for a barber and a tailor.
\end{quote}

\subsection*{Your turn}
Say these aloud:

\begin{itemize}
\raggedright
  \item cotton, silk, leather, brass, lime
  \item manure, cow dung, a weed, elephant-foot yam, a pumpkin
  \item a coin, a banknote, small change, the balance, a discount
\end{itemize}

\begin{tcolorbox}[breakable,colback=teal!4,colframe=teal!35!black,title={Before you move on}]
Cotton? (\textbf{\ml{പരുത്തി}}, \emph{parutti}.) Silk? (\textbf{\ml{പട്ട്}}, \emph{paṭṭŭ}.) Leather? (\textbf{\ml{തുകൽ}}, \emph{tukal}.) Brass? (\textbf{\ml{പിച്ചള}}, \emph{piccaḷa}.) Lime? (\textbf{\ml{ചുണ്ണാമ്പ്}}, \emph{cuṇṇāmpŭ}.) Manure? (\textbf{\ml{വളം}}, \emph{vaḷaṁ}.) Cow dung? (\textbf{\ml{ചാണകം}}, \emph{cāṇakaṁ}.) A weed? (\textbf{\ml{കള}}, \emph{kaḷa}.) Elephant-foot yam? (\textbf{\ml{ചേന}}, \emph{cēna}.) A pumpkin? (\textbf{\ml{മത്തങ്ങ}}, \emph{mattaṅṅa}.) A coin? (\textbf{\ml{നാണയം}}, \emph{nāṇayaṁ}.) A banknote? (\textbf{\ml{നോട്ട്}}, \emph{nōṭṭŭ}.) Small change? (\textbf{\ml{ചില്ലറ}}, \emph{cillaṟa}.) The balance? (\textbf{\ml{ബാക്കി}}, \emph{bākki}.) A discount? (\textbf{\ml{കിഴിവ്}}, \emph{kiḻivŭ}.)
\end{tcolorbox}
\section[(dialogue)]{Words from chapters 220-222 — a second pass}

\label{lesson:ML-R222-second-pass-words-from-chapters-220-222}



\begin{quote}
Say the words for a colleague and an official.
\end{quote}

\subsection*{Your turn}
Say these aloud:

\begin{itemize}
\raggedright
  \item a colleague, an official, a clerk, an engineer, a lawyer
  \item a word, a sentence, fees, a classmate, education
  \item geography, literature, grammar, a barber, a tailor
\end{itemize}

\begin{tcolorbox}[breakable,colback=teal!4,colframe=teal!35!black,title={Before you move on}]
A colleague? (\textbf{\ml{സഹപ്രവർത്തകൻ}}, \emph{sahapravarttakan}.) An official? (\textbf{\ml{ഉദ്യോഗസ്ഥൻ}}, \emph{udyōgasthan}.) A clerk? (\textbf{\ml{ഗുമസ്തൻ}}, \emph{gumastan}.) An engineer? (\textbf{\ml{എഞ്ചിനീയർ}}, \emph{eñcinīyar}.) A lawyer? (\textbf{\ml{വക്കീൽ}}, \emph{vakkīl}.) A word? (\textbf{\ml{വാക്ക്}}, \emph{vākkŭ}.) A sentence? (\textbf{\ml{വാക്യം}}, \emph{vākyaṁ}.) Fees? (\textbf{\ml{ഫീസ്}}, \emph{phīsŭ}.) A classmate? (\textbf{\ml{സഹപാഠി}}, \emph{sahapāṭhi}.) Education? (\textbf{\ml{വിദ്യാഭ്യാസം}}, \emph{vidyābhyāsaṁ}.) Geography? (\textbf{\ml{ഭൂമിശാസ്ത്രം}}, \emph{bhūmiśāstraṁ}.) Literature? (\textbf{\ml{സാഹിത്യം}}, \emph{sāhityaṁ}.) Grammar? (\textbf{\ml{വ്യാകരണം}}, \emph{vyākaraṇaṁ}.) A barber? (\textbf{\ml{ബാർബർ}}, \emph{bārbar}.) A tailor? (\textbf{\ml{തയ്യൽക്കാരൻ}}, \emph{tayyalkkāran}.)
\end{tcolorbox}
